\section{Συμπεράσματα και Μελλοντικές Επεκτάσεις}
\label{ConclusionsFuture}
\subsection{Συμπεράσματα}
\label{Conclusions}
Από τη συγκεκριμένη διπλωματική εργασία, εξήχθησαν διάφορα συμπεράσματα τα οποία αφορούν τόσο τη φύση του προβλήματος της αναγνώρισης ομιλίας από τα χείλη (lip reading) όσο και τη συγκεκριμένη μεθοδολογία που χρησιμοποιήθηκε και δοκιμάστηκε (τόσο του αρχικού visual μοντέλου, βλ. \ref{ModelStructure}, \ref{Experiment1},  όσο και του audiovisual μοντέλου, βλ. \ref{AudiovisualModel}, \ref{Experiment2}). Τα σημαντικά από αυτά συμπεριλαμβάνονται παρακάτω.

Γενικά, για το πρόβλημα αναγνώρισης φωνημάτων:

\begin{itemize}
\item Το πρόβλημα πρόβλεψης φωνημάτων (και, γενικότερα, ανθρώπινης ομιλίας) χρησιμοποιώντας μόνο οπτικά στοιχεία (βουβό βίντεο) είναι αδύνατον να επιλυθεί πλήρως, λόγω της ύπαρξης των visemes, των φωνημάτων δηλαδή που, ενώ είναι διαφορετικά, παρουσιάζουν την ίδια εικόνα στα χείλη κατά την εκφορά τους.

\item Για να μειωθεί η επίδραση των visemes στην αναγνώριση των φωνημάτων, είναι σημαντικό να χρησιμοποιείται κατά την πρόβλεψη παράθυρο μεγαλύτερο από τη διάρκεια του ίδιου του φωνήματος, έτσι ώστε το εκάστοτε μοντέλο πρόβλεψης να μπορεί να αξιοποιεί στοιχεία από τη γειτονιά του ίδιου φωνήματος (ποια προηγούνται και ποια έπονται) και, είτε επειδή χτίζει γνώση του λεξιλογίου σε περιορισμένο περιβάλλον, είτε επειδή μαθαίνει τις γενικές τάσεις κάποιας γλώσσας όσον αφορά την αλληλουχία φωνημάτων σε μη περιορισμένο περιβάλλον, να μαντεύει ορθά φωνήματα ακόμα και όταν αυτά ανήκουν σε κάποιο viseme.

\item Η κατάτμηση του προβλήματος της αναγνώρισης της ομιλίας πρώτα στην αναγνώριση φωνημάτων και έπειτα στον συνδυασμό των προβλεπόμενων φωνημάτων για τη δημιουργία λέξεων φαίνεται να είναι ουσιώδης, διότι η αναγνώριση ολόκληρης φράσης δίχως το ενδιάμεσο βήμα της πρόβλεψης των φωνημάτων θα ήταν πιθανότατα διαδικασία ιδιαίτερα επιρρεπής στον θόρυβο και στις ιδιομορφίες της προφοράς του κάθε ανθρώπου.
\end{itemize}

Όσον αφορά τη δομή του προτεινόμενου visual μοντέλου και τα πειραματικά αποτελέσματα που παρουσιάζει:

\begin{itemize}
\item Από το πείραμα αφαίρεσης που παρουσιάστηκε στην ενότητα \ref{Experiment1}, συμπεραίνουμε ότι απαιτείται ο συνδυασμός του Feature Extractor, ο οποίος στη συγκεκριμένη περίπτωση είναι ένα συνελικτικό δίκτυο, και του Sequence Encoder, ο οποίος είναι υπεύθυνος για τον συνδυασμό των features που εξήχθησαν και την παραγωγή της πληροφορίας από την οποία τελικά γίνεται η πρόβλεψη φωνημάτων. Το μοντέλο δίχως τον Sequence Encoder παρουσίασε τη χειρότερη επίδοση από αυτά που δοκιμάστηκαν.

\item Συγκεκριμένα για τον Sequence Encoder, από τη μελέτη αφαίρεσης καταλαβαίνουμε ότι το BiLSTM υπερέχει του Transformer, κάτι το οποίο συμβαίνει πιθανότατα λόγω του μικρού μεγέθους του χρησιμοποιούμενου dataset (GRID) και όχι λόγω κάποιας σχεδιαστικής υπεροχής του μοντέλου με BiLSTM. Όταν τα δεδομένα δεν είναι πολλά, οι transformers τείνουν να υπερπροσαρμόζονται (overfit) στα δεδομένα εκπαίδευσης και τελικά να έχουν χειρότερη επίδοση από άλλες παραδοσιακές αρχιτεκτονικές. Σε περίπτωση δοκιμής σε κάποιο μεγαλύτερο, ειδικά in the wild dataset, η εικασία είναι ότι ο Transformer θα υπερείχε του BiLSTM.

\item Όσον αφορά την αντικειμενική συνάρτηση του μοντέλου, συμπεραίνουμε ότι ο συνδυασμός επίβλεψης τόσο σε επίπεδο frame όσο και σε επίπεδο τελικής αλληλουχίας των προβλεπόμενων φωνημάτων υπερέχει της χρήσης οποιασδήποτε από τις δύο μεθόδους ξεχωριστά. Ο όρος CTC της αντικειμενικής συνάρτησης είναι υπεύθυνος για τη πρόβλεψη αλληλουχιών φωνημάτων που έχουν νόημα, ενώ ο όρος cross-entropy προσφέρει σήμα που εκπαιδεύει το μοντέλο στους χρονισμούς εκφοράς του κάθε φωνήματος ξεχωριστά (σημαντικό, αν στόχος είναι η αναπαραγωγή της ομιλίας και όχι αποκλειστικά η πρόβλεψή της).

\item Η χρήση soft ή hard targets κατά την εκπαίδευση φαίνεται να παρουσιάζει, τουλάχιστον στο περιορισμένο περιβάλλον, τη μικρότερη επίδραση στο τελικό αποτέλεσμα. Η επίδραση της επιλογής soft targets έναντι hard targets ίσως αυξάνεται με το μέγεθος και την ποικιλομορφία του dataset, καθώς σε περιορισμένο περιβάλλον είναι πιο εύκολο να ``αποστηθιστούν'' συγκεκριμένες αλληλουχίες φωνημάτων, μια και αυτές που προκύπτουν χρησιμοποιώντας οπτικά παρόμοια φωνήματα είναι τελικά άκυρες.

\item Η χρήση της depthwise separable συνέλιξης έναντι της παραδοσιακής εκδοχής της παρέχει σημαντική μείωση στις εκπαιδευόμενες παραμέτρους του μοντέλου, αυξάνοντας την ταχύτητα εκπαίδευσης και μειώνοντας το υπολογιστικό κόστος.

\item Το μοντέλο που αναδεικνύεται βέλτιστο από τη μελέτη αφαίρεσης (αυτό που τελικά προτείνεται στην ενότητα \ref{ModelStructure}) παρουσιάζει ευχάριστα αποτελέσματα στην πρόβλεψη των φωνημάτων, εικόνα που γίνεται ιδιαίτερα εμφανής μέσω του πίνακα σύγχυσης. Η πρόβλεψη των εκάστοτε λέξεων στο περιορισμένο περιβάλλον από τα προβλεπόμενα φωνήματα φαίνεται επίσης να είναι ιδιαίτερα επιτυχής, με μοναδική εξαίρεση την κατηγορία των γραμμάτων στο GRID. Η αδυναμία αυτή οφείλεται στον αριθμό των υποψήφιων γραμμάτων, ο οποίος είναι πολύ μεγαλύτερος από τους αντίστοιχους αριθμούς για τις άλλες κατηγορίες, αλλά και στο μέγεθος των γραμμάτων σε φωνήματα. Καθώς κατά την εκφορά των περισσότερων γραμμάτων χρησιμοποιούνται περίπου 2 φωνήματα, η διαδικασία πρόβλεψης είναι ιδιαίτερα ευάλωτη στον θόρυβο που προκύπτει από τη λανθασμένη πρόβλεψη ενός από αυτά. Ο προβλεπόμενος χρονισμός του κάθε φωνήματος (και τελικά κάθε λέξης της φράσης) φαίνεται να είναι ιδιαίτερα ακριβής.
\end{itemize}

Όσον αφορά τη δομή του προτεινόμενου audiovisual μοντέλου και τα πειραματικά αποτελέσματα που παρουσιάζει:

\begin{itemize}
\item Η ύπαρξη του ηχητικού σήματος, όπως ήταν αναμενόμενο, βοηθά σε πολύ σημαντικό βαθμό την πρόβλεψη των φωνημάτων και τελικά της ομιλίας, ακόμα και όταν το σήμα αυτό είναι θορυβώδες. Το ηχητικό σήμα προσφέρει τη λύση στο πρόβλημα των visemes, διότι τα visemes ακούγονται διαφορετικά, ακόμα και αν παράγονται τοποθετώντας τα χείλη με παρόμοιο τρόπο.

\item Η χρήση των Mel-Spectrograms και όχι του ίδιου του ηχητικού σήματος ως είσοδο στο audio stream του μοντέλου εξυπηρετεί τη διατήρηση της γενικής ιδέας του visual μοντέλου, συγκεκριμένα τη χρήση συνελικτικού νευρωνικού δικτύου για την εξαγωγή των features. Επιπροσθέτως, συμβάλλει στη μείωση του όγκου των δεδομένων που τελικά εισέρχονται στο μοντέλο.

\item Σε περιπτώσεις ύπαρξης θορύβου στο ηχητικό σήμα αλλά και στο βίντεο, είναι σημαντικό για ένα μοντέλο που αξιοποιεί τα δύο αυτά σήματα να μάθει πότε πρέπει να σέβεται την πρόβλεψη που εξάγεται από τη μία τροπικότητα έναντι της άλλης. Επίσης, πρέπει να μάθει να συνδυάζει τις πληροφορίες των δύο τροπικοτήτων ώστε να εξάγει ένα κοινό συμπέρασμα για τα φωνητικά χαρακτηριστικά της εκάστοτε χρονικής στιγμής.

\item Το πείραμα 2 της ενότητας \ref{Experiment2} αποδεικνύει την ικανότητα του μοντέλου να εμπιστεύεται μία τροπικότητα όλο και λιγότερο, όσο αυτή γίνεται λιγότερο αξιόπιστη (και το βίντεο, αλλά και τον ήχο). Αποδεικνύεται, επίσης, η ανώτερη ποιότητα της πρόβλεψης όταν συνυπάρχουν οι δύο τροπικότητες έναντι της αξιοποίησης μονάχα μίας εκ των δύο. Η ποιότητα της πρόβλεψης, όπως είναι αναμενόμενο, μειώνεται όσο μειώνεται η αξιοπιστία οποιασδήποτε από τις δύο τροπικότητες.

\item Το audiovisual μοντέλο φαίνεται να εμφανίζει μία ελαφριά προτίμηση για τα features που προέρχονται από το ηχητικό σήμα, κρίνοντας από τις τιμές της audio gate και τις αντίστοιχες της video gate. Ακόμα και όταν ο θόρυβος είναι συνεχής για το ηχητικό σήμα αλλά πολύ σύντομος για το βίντεο, το μοντέλο εμπιστεύεται ελάχιστα περισσότερο τον ήχο. Αυτό οφείλεται πιθανότατα στο γεγονός ότι η επίλυση του προβλήματος πρόβλεψης φωνημάτων είναι πολύ πιο εύκολη αξιοποιώντας την ομιλία σε σχέση με το βίντεο της ομιλίας.
\end{itemize}

\subsection{Μελλοντικές επεκτάσεις}
\label{Future}

Παρακάτω γίνεται αναφορά σε πιθανές μελλοντικές επεκτάσεις της συγκεκριμένης εργασίας.

\begin{itemize}

\item Η βασική μελλοντική επέκταση της εργασίας αυτής είναι η δοκιμή της προτεινόμενης μεθοδολογίας σε περιβάλλον το οποίο είναι μη περιορισμένο (in the wild), για παράδειγμα στα ευρέως διαδεδομένα datasets LRS2 και LRS3 \cite{Afouras18c}, \cite{Chung16}, \cite{Chung17}, \cite{Chung17a}, στα οποία δεν κατέστη δυνατή η απόκτηση πρόσβασης κατά τη διάρκεια της εργασίας. Το πείραμα αυτό κρίνεται ιδιαίτερα σημαντικό, τόσο επειδή ανταποκρίνεται σε πραγματικές συνθήκες στις οποίες μπορεί να εφαρμοστεί η αναγνώριση της ανθρώπινης ομιλίας (προφανώς η αναγνώριση σε in the wild περιβάλλον είναι πολύ πιο δύσκολο πρόβλημα και ταυτόχρονα πολύ πιο χρήσιμη εφαρμογή από την αναγνώριση σε περιορισμένο περιβάλλον με μικρό λεξιλόγιο), αλλά και επειδή πρόκειται για το benchmark test των μοντέλων του state-of-the-art. Είναι σημαντικό να μελετηθεί κατά πόσο μειώνεται η επίδοση των προβλέψεων φωνημάτων του μοντέλου στο in the wild dataset σε σχέση με το περιορισμένο dataset και ποιες τεχνικές μπορούν να εφαρμοστούν για τη μείωση του φαινομένου αυτού (π.χ. διαφορετικές μορφές data augmentation που εστιάζουν σε σπάνιες αλληλουχίες φωνημάτων στις οποίες δυσκολεύεται το μοντέλο). Κρίνεται, επίσης, ενδιαφέρουσα η μελέτη για το αν ο Transformer όντως ανταπεξέρχεται καλύτερα ως Sequence Encoder στο μεγαλύτερο και πιο ποικιλόμορφο dataset της in the wild περίπτωσης σε σχέση με το BiLSTM.

\item Για την επίτευξη των στόχων του πειράματος αυτού, απαιτείται και η ανάπτυξη μεθόδου πρόβλεψης λέξεων και, τελικά, φράσεων με δεδομένο την ύπαρξη σε μη περιορισμένο περιβάλλον. Όπως παρουσιάζεται στην ενότητα \ref{PhrasePrediction}, μπορεί να δοκιμαστεί LLM με κατάλληλο fine-tuning ώστε να προβλέπει μία πιθανή φράση από την αλληλουχία των προβλεπόμενων φωνημάτων. Προφανώς, λόγω του τεράστιου αριθμού των ποσοτήτων, η αλγοριθμική πρόβλεψη δεν μπορεί να εφαρμοστεί στην περίπτωση του in the wild περιβάλλοντος.

\item Όσον αφορά την εφαρμογή της πρόβλεψης της ομιλίας, μία σημαντική μελλοντική επέκταση είναι η αναγνώριση λέξεων σε πραγματικό χρόνο. Με την υπάρχουσα εκδοχή, το μοντέλο απαιτεί τη ροή των ήδη καταγεγραμμένων δεδομένων (είτε αποκλειστικά βίντεο είτε βίντεο και ήχου). Σε εφαρμογές Augmented Reality, κατά τις οποίες η πρόβλεψη των λέξεων εφαρμόζεται σε πραγματικό χρόνο ώστε να βοηθήσει είτε άτομα με ειδικές ανάγκες (π.χ. προβλήματα ακοής) είτε σε περιπτώσεις ύπαρξης θορύβου ο οποίος καθιστά την ομιλία μη κατανοητή (π.χ. θορυβώδες αμφιθέατρο), τα δεδομένα εισέρχονται στο μοντέλο την ίδια στιγμή που παράγονται. Έτσι, κρίνεται ενδιαφέρουσα μελλοντική επέκταση η κατάλληλη τροποποίηση του μοντέλου ώστε να δέχεται σε πραγματικό χρόνο τα δεδομένα και να αποφασίζει για τις λέξεις που εκφέρονται.

\item Μία ακόμη πιθανή γενίκευση είναι η υποστήριξη διαφόρων γλωσσών και όχι μόνο της αγγλικής. Προφανώς το πρόβλημα αυτό είναι πολύ πιο δύσκολο, καθώς προϋποθέτει το μοντέλο να μπορεί να αντιληφθεί τη γλώσσα που ομιλείται ανά πάσα στιγμή, είτε αποκλειστικά από το βίντεο είτε αξιοποιώντας θορυβώδη ήχο και βίντεο. Ακόμη, θα ήταν αναγκαία η απόκτηση ενός πολύ γενικευμένου dataset που περιλαμβάνει βίντεο από ομιλίες διαφόρων γλωσσών του κόσμου, αυξάνοντας το υπολογιστικό κόστος σημαντικά.

\item Όσον αφορά τον εξαγωγέα χαρακτηριστικών του βίντεο, ενδιαφέρουσα επέκταση είναι η συγκριτική μελέτη της depthwise separable 3D συνέλιξης, η οποία χρησιμοποιείται στην εργασία, με την πλήρη 3D συνέλιξη (βλ. ενότητα \ref{ModelStructure}). Η depthwise separable συνέλιξη μειώνει σημαντικά τον αριθμό των παραμέτρων και το υπολογιστικό κόστος, αλλά ενισχύει το inductive bias του μοντέλου, καθώς η μίξη των καναλιών δεν μπορεί να εξαρτάται από το χωροχρονικό παρακείμενο εντός ενός block. Μία μελέτη αφαίρεσης μεταξύ των δύο επιλογών θα έδειχνε κατά πόσο η απώλεια εκφραστικότητας επηρεάζει την ποιότητα της μάθησης, όπως και αν η επίδραση αυτή μεταβάλλεται με το μέγεθος του dataset (π.χ. σε ένα μεγαλύτερο, in the wild dataset). Αντίστοιχη σύγκριση μπορεί να γίνει και για τον εξαγωγέα χαρακτηριστικών του ήχου στο audiovisual μοντέλο, ο οποίος χρησιμοποιεί depthwise separable 2D συνέλιξη.

\item Όσον αφορά το audiovisual μοντέλο, ενδιαφέρουσα επέκταση είναι η συστηματική μελέτη της εισόδου του σταδίου Reliability-gated Fusion (βλ. ενότητα \ref{ReliabilityFusion}). Στην παρούσα εργασία, πέραν των αναπαραστάσεων $F_v\sp{\prime}(t)$ και $F_a\sp{\prime}(t)$ και των σημάτων αξιοπιστίας, η είσοδος περιλαμβάνει την απόλυτη διαφορά ανά στοιχείο $\lvert F_v\sp{\prime}(t) - F_a\sp{\prime}(t) \rvert$, χωρίς να έχει μελετηθεί η επίδρασή της. Θα μπορούσαν να συγκριθούν, μέσω κατάλληλης μελέτης αφαίρεσης, εκδοχές στις οποίες η διαφορά δε χρησιμοποιείται καθόλου, χρησιμοποιείται μόνο το μέτρο (η νόρμα) της διαφοράς $\lVert F_v\sp{\prime}(t) - F_a\sp{\prime}(t) \rVert$ ως ένα βαθμωτό μέγεθος, ή χρησιμοποιείται ολόκληρο το διάνυσμα της διαφοράς $F_v\sp{\prime}(t) - F_a\sp{\prime}(t)$ (διατηρώντας και το πρόσημο), ώστε να διαπιστωθεί κατά πόσο η πληροφορία αυτή βοηθά πράγματι το μοντέλο να αναγνωρίζει την τροπικότητα που επηρεάζεται από τον θόρυβο.

\end{itemize}